\section{Trace, Report, and Evidence Model}\label{sec:evidence}

Each complete request directory contains five artifacts:
\code{session.argx}, compiled \code{session.argbc.json},
\code{session.trace.json}, \code{session.security.json}, and
\code{session.evidence.json}.  The trace is an ordered ledger of compiler/VM
and policy events.  The SecurityReport is a structured interpretation of
runtime evidence, including execution status, policy results, review state,
provider boundary, and declared controls.  The EvidenceBundle links source,
bytecode, trace, report, and ledger through semantic SHA-256 digests.

\begin{figure}[t]
\centering
\includegraphics[width=\linewidth]{artifact-schema.pdf}
\caption{Schema-level view of the five-artifact request directory. The evidence
bundle references normalized content digests; it does not contain a security
verdict.}
\label{fig:artifact-schema}
\end{figure}

A semantic digest is computed over canonical JSON derived from a deserialized
value, avoiding dependence on whitespace or object-key formatting.  Bundle
paths resolve relative to the bundle directory.  Offline verification loads
the referenced artifacts, recomputes their semantic digests, and compares them
with the bundle.  It requires neither network access nor provider credentials.
This verifies internal linkage: the inspected artifacts correspond to the
values named by the bundle.

Three distinctions are essential.  First, a valid digest detects content
change relative to the bundle but does not prove who created the bundle.
Second, a linked trace does not prove that events outside the instrumented
runtime did not occur.  Third, artifact integrity does not turn a failing
policy result into approval.  No signature, trusted timestamp, transparency
service, hardware root, or independent witness is evaluated here.

The local trust ledger is similarly bounded.  Hash-linked entries and validated
references can expose accidental or adversarial modification when the verifier
has an expected root or bundle.  They do not provide distributed consensus,
blockchain immutability, or non-repudiation.  The declarations explicitly deny
such security claims.

\begin{figure}[t]
\centering
\includegraphics[width=\linewidth]{claim-boundaries.pdf}
\caption{Claim-boundary plate separating implemented mechanisms, declarative
metadata, proposed architecture, and properties not claimed.}
\label{fig:claim-boundaries}
\end{figure}

\begin{table}[t]
\centering
\caption{Interpretive boundaries applied to every result in this study.}
\label{tab:claim-boundaries}
\begin{tabular}{@{}lllll@{}}
\toprule
Concept & Implemented & Declarative & Proposed & Not claimed \\
\midrule
Provider boundary & simulated allowlist & contract metadata & -- & external execution \\
Trust evidence & digest bundle & ATrust map & -- & live attestation \\
Discovery & local catalog & sovereign metadata & federation & operational DNS \\
Security scope & fail-closed controls & threat mappings & deployment study & certification \\
\bottomrule
\end{tabular}

\end{table}
